\chapter[Exact Solution of the Ising Model]{Exact Solution of the \\Ising Model}

 [...]

\section{Ising Model in One Dimension}

The hamiltonian of the 1D Ising model is given by
\[
    H = -J \sum_{i=1}^{N} s_i s_{i+1} - h \sum_{i=1}^{N} s_i,
\]
where we have periodic boundary conditions \(s_{N+1} = s_1\). The partition function depends upon the temperature \(T\), the external magnetic field \(h\) and the number of spins \(N\), and it can be written as
\[
    Z_N(\beta, h) = \sum_{\{s_i\}} e^{-\beta H(\{s_i\})}.
\]
It is important since we can derive an expression for the ensemble average of the magnetization per site as a function of the partition function:
\[
    \mathbb{E}_c \left[ \frac{1}{N} \sum_{i=1}^{N} s_i \right] = - \frac{1}{N\beta} \frac{\partial \log Z_N}{\partial h}.
\]
If we manipulate the expression of the hamiltonian, to express two symmetric sums, we get
\[
    H = -J \sum_{i=1}^{N} s_i s_{i+1} - \frac{h}{2} \sum_{i=1}^{N} s_i s_{i+1}
\]
\[
    Z = \sum_{\{s_i\}} \prod_{i=1}^{N} (e^{\beta J s_i s_{i+1} + \frac{\beta h}{2} s_i s_{i+1}})\dots
\]

Now we can introduce the \textbf{transfer matrix} \(T\), which is a \(2 \times 2\) matrix that encodes the interactions between neighboring spins and the effect of the external magnetic field. The transfer matrix is defined as
\[
    T_{s_i, s_{i+1}} = e^{\beta J s_i s_{i+1} + \frac{\beta h}{2} (s_i + s_{i+1})},
\]
where \(s_i\) and \(s_{i+1}\) can take values of \(\pm 1\). The partition function can then be expressed as a product of transfer matrices:
\[
    Z = \sum_{\{s_i\}} T_{s_1, s_2} T_{s_2, s_3} \cdots T_{s_N, s_1} = \text{Tr}(T^N),
\]
where the trace is taken over the possible spin states. The transfer matrix method allows us to compute the partition function and other thermodynamic quantities by diagonalizing the transfer matrix and finding its eigenvalues. The largest eigenvalue of the transfer matrix dominates the behavior of the partition function in the thermodynamic limit, and it can be used to derive expressions for the free energy, magnetization, and other physical quantities of interest. This method is particularly powerful for solving one-dimensional models like the 1D Ising model, where it provides an exact solution for the partition function and allows us to analyze the behavior of the system at different temperatures and external fields.

\subsubsection{Transfer Matrix}

Expressing the partition function in terms of the transfer matrix, we can notice how the computational complexity of the problem is reduced from summing over \(2^N\) configurations to diagonalizing a \(2 \times 2\) matrix multiplied \(N\) times with itself, which is a much more tractable problem. The eigenvalues of the transfer matrix can be computed explicitly, and they determine the thermodynamic properties of the system. We can write the transfer matrix explicitly as
\[
    T = \begin{pmatrix}
        e^{\beta J + \beta h} & e^{-\beta J}          \\
        e^{-\beta J}          & e^{\beta J - \beta h}
    \end{pmatrix},
\]
with the idea of diagonalizing this matrix to find its eigenvalues \(\lambda_1\) and \(\lambda_2\). Let us note that \(T = T^{\dagger}\), thus we know there exist an orthogonal matrix \(O = O^T\) such that \(T = O D O^{T}\), where \(D\) is a diagonal matrix containing the eigenvalues of \(T\). The diagonal maatrix \(D\) can be expressed as
\[
    D = \begin{pmatrix}
        \lambda_+ & 0         \\
        0         & \lambda_-
    \end{pmatrix},
\]
and we can compute the the \(N\)th power of the transfer matrix as
\[
    T^N = (ODO^T)^N = (ODO^T)(ODO^T)\cdots(ODO^T) = O D^N O^T,
\]
where we have used the fact that \(O^T O = I\) to simplify the expression. The partition function can then be expressed as
\[
    Z = \text{Tr}(T^N) = \text{Tr}(O D^N O^T) = \text{Tr}(D^N) = \lambda_+^N + \lambda_-^N,
\]
where \(\lambda_+\) and \(\lambda_-\) are the eigenvalues of the transfer matrix \(T\) and we have used the cyclic property of the trace.

Now we can also compute the determinant of the transfer matrix, which is given by
\[
    \det(T) = \lambda_+ \lambda_- = e^{2\beta J} - e^{-2\beta J} = 2 \sinh(2\beta J),
\]
while for the trace we had
\[
    \text{Tr}(T) = \lambda_+ + \lambda_- = e^{\beta J + \beta h} + e^{\beta J - \beta h} + 2 e^{-\beta J} = 2 e^{\beta J} \cosh(\beta h) + 2 e^{-\beta J}.
\]
The eigenvalues of the transfer matrix can be computed using the characteristic polynomial, which is given by
\[
    x^2 - \text{Tr}(T) x + \det(T) = 0,
\]
which, inserting the expressions for the trace and determinant, can be rewritten as
\[
    \dots
\]
and in the end we get the eigenvalues as
\[
    \lambda_{\pm} = e^{\beta J} \left(\cosh(\beta h) \pm \sqrt{\sinh^2(\beta h) + e^{-4\beta J}}\right),
\]
which gives
\[
    \vert \frac{\lambda_-}{\lambda_+} \vert <1 , \quad \vert \frac{\lambda_-}{\lambda_+} \vert^N \ll 1.
\]
thus we can write the partition function in the thermodynamic limit as
\[
    Z = \lambda_+^N + \lambda_-^N = \lambda_+^N \left(1 + \left(\frac{\lambda_-}{\lambda_+}\right)^N\right) \approx \lambda_+^N,
\]
which allows us to compute the free energy per site as
\[
    f = -\frac{1}{\beta} \lim_{N \to \infty} \frac{1}{N} \log Z,
\]
where \(f = \lim_{N \to \infty} \frac{F}{N}\); thus we have
\[
    f = -\frac{1}{\beta} \log \lambda_+ = -J - \frac{1}{\beta} \log\left( \cosh(\beta h) + \sqrt{\sinh^2(\beta h) + e^{-4\beta J}} \right).
\]
from this we can understand that there are no singularities in the free energy as a function of temperature \(T\) and external field \(h\), which indicates that there is no phase transition in the 1D Ising model at finite temperature. This is consistent with the known result that the 1D Ising model does not exhibit a phase transition at finite temperature, and it highlights the importance of dimensionality and fluctuations in determining the behavior of the system near critical points.

\subsection{Magnetization}

Let us now compute the magnetization per site as a function of temperature \(T\) and external field \(h\). We can use the expression for the partition function to compute the magnetization as
\[
    \mathbb{E}_{c}\left[s_i\right] = m = -\frac{1}{N} \frac{\partial}{\partial h} \log Z = -\frac{1}{\beta} \frac{\partial}{\partial h} (\frac{F}{N}).
\]
Inserting the expression for the free energy per site, we get
\[
    m = \frac{\sinh(\beta h)}{\sqrt{\sinh^2(\beta h) + e^{-4\beta J}}}.
\]
This expression shows that the magnetization \(m\) is a smooth function of temperature \(T\) and external field \(h\), and it does not exhibit any singularities or discontinuities, which is consistent with the absence of a phase transition in the 1D Ising model at finite temperature.

\TODO{add plot of m vs h for different T.}

As \(T\) increases, the curve becomes more steep, but it never exihibit a discontinuity or a singularity, which confirms that there is no phase transition in the 1D Ising model at finite temperature. Only in the limit \(T \to 0\) we have a discontinuity in the magnetization at \(h = 0\), where the system transitions from a fully magnetized state with \(m = +1\) for \(h > 0\) to a fully magnetized state with \(m = -1\) for \(h < 0\):
\[
    \lim_{T \to 0} m(T,h) = \text{sgn}(\sinh(\beta h)) = \text{sgn}(h).
\]

We could repeat the computation for \(\chi\) the magnetic susceptibility, which is given by
\[
    \chi = \frac{\partial}{\partial h} m(T,h) = \dots
\]
and it can be shown that \(\chi\) does not diverge at any finite temperature, which is consistent with the absence of a phase transition in the 1D Ising model at finite temperature. Only in the limit \(T \to 0\) we have a divergence in the magnetic susceptibility at \(h = 0\), which indicates a critical point at zero temperature.

The euristic derivation of mean field theory has to have some assumptions which do not hold at low dimensions, and this is the reason why it fails to capture the correct critical behavior in 1D and 2D systems. Only in higher dimensions, where fluctuations are less important, the mean field theory provides a more accurate description of the critical behavior and the critical exponents.

The exact solution of the 1D Ising model using the transfer matrix method allows us to understand the limitations of mean field theory and the importance of fluctuations in determining the behavior of the system near critical points.

\subsection{Two Point Correlation Function}

It is now time to compute the two-point correlation function for the 1D Ising model: the correlation function is defined as
\[
    \mathbb{E}_{c}\left[ s_i s_{i+r} \right] = \frac{1}{Z} \sum_{\{s_k\}} s_i s_{i+r} e^{-\beta H(\{s_k\})},
\]
which can become a connected correlation function if we subtract the product of the magnetizations:
\[
    \mathbb{E}_{c}\left[ s_i s_{i+r} \right] - \mathbb{E}_{c}\left[ s_i \right] \mathbb{E}_{c}\left[ s_{i+r} \right] = \frac{1}{Z} \sum_{\{s_k\}} s_i s_{i+r} e^{-\beta H(\{s_k\})} - m^2.
\]
To compute the correlation function, we can use the transfer matrix method again. We can express the correlation function in terms of the transfer matrix as follows:
\[
    \mathbb{E}_{c}\left[ s_i s_{i+r} \right] = \frac{1}{Z} \sum_{\alpha_1=1}^2 \sum_{\alpha_2=1}^2 \dots \sum_{\alpha_N=1}^2 \sigma_{\alpha_j} \sigma_{\alpha_{j+r}} T_{\alpha_1, \alpha_2} T_{\alpha_2, \alpha_3} \cdots T_{\alpha_N, \alpha_1},
\]
where \(\sigma_{\alpha_j}\) is the spin value corresponding to the state \(\alpha_j\) (i.e., \(\sigma_1 = +1\) and \(\sigma_2 = -1\)). We can rewrite this expression in a more compact form using matrix notation. Let us define a diagonal matrix \(D\) such that
\[
    D = \begin{pmatrix}
        1 & 0  \\
        0 & -1
    \end{pmatrix}, \text{ and } D_{\alpha,\beta} = \delta_{\alpha,\beta} \cdot \sigma_{\alpha},
\]
meaning that \(D\) is a diagonal matrix with entries corresponding to the spin values. Then we can express the correlation function as
\[
    \mathbb{E}_{c}\left[ s_j s_{j+r} \right] = \frac{1}{Z} \sum_{\alpha_1=1}^2 \sum_{\alpha_2=1}^2 \dots \sum_{\alpha_N=1}^2 T_{\alpha_1, \alpha_2} T_{\alpha_2, \alpha_3} \cdots T_{\alpha_{j-1}, \alpha_j} D_{\alpha_j, \beta_1} T_{\beta_1, \alpha_{j+1}} \cdots T_{\alpha_{j+r-1}, \alpha_{j+r}} D_{\alpha_{j+r}, \beta_2} T_{\beta_2, \alpha_{j+r+1}} \cdots T_{\alpha_N, \alpha_1},
\]
where we have inserted the diagonal matrices \(D\) at the positions corresponding to the spins \(s_j\) and \(s_{j+r}\). This expression can be further simplified by recognizing that the sums over the indices \(\alpha_k\) correspond to matrix multiplications, and we can express
\[
    \mathbb{E}_{c}\left[ s_j s_{j+r} \right] = \frac{1}{Z} \text{Tr}(T^{j-1} D T^{r} D T^{N-r+1-j}),
\]
where we have used the cyclic property of the trace to rearrange the order of the matrices. This expression allows us to compute the two-point correlation function in terms of the transfer matrix and the diagonal matrix \(D\), which encodes the spin values. By diagonalizing the transfer matrix with a unitary matrix \(U\):
\[
    T = U \begin{pmatrix}
        \lambda_+ & 0         \\
        0         & \lambda_-
    \end{pmatrix} U^{\dagger},
\]
where \(U\) can be computed explicitly, we can express the correlation function in terms of the eigenvalues and eigenvectors of the transfer matrix, which allows us to analyze its behavior as a function of the distance \(r\) and the temperature \(T\). Let us focus on the case without external field: the unitary transformation becomes
\[
    \dots
\]
and so we can write the transfer matrix in the diagonal basis as
\[
    \dots
\]
Now putting everything together, we can express the correlation function as
\[
    \begin{aligned}
        \dots
    \end{aligned}
\]
In the end we find a close analytical form for the correlation function, which can be expressed as
\[
    \mathbb{E}_{c}\left[ s_j s_{j+r} \right]_{h=0} = \frac{\lambda_+^{N-r}\lambda_-^r + \lambda_-^{N-r}\lambda_+^r}{\lambda_+^N + \lambda_-^N},
\]
and again if we use tha fact that \(\vert \frac{\lambda_-}{\lambda_+} \vert < 1\), we can simplify this expression in the thermodynamic limit as
\[
    \mathbb{E}_{c}\left[ s_j s_{j+r} \right]_{h=0} =\dots \sim e^{-\tfrac{r}{\xi}},
\]
effectively defining the correlation length \(\xi\) as
\[
    \xi^{-1} = \log(\frac{\lambda_+}{\lambda_-}) > 0,
\]
which is a measure of how far correlations between spins extend in the system.
Now considering the absence of an external field, we can express the eigenvalues of the transfer matrix as
...

and in the end we get the correlation length in the \(T \to 0\) limit without external field as
\[
    \xi \sim e^{2\beta J},
\]
so that it diveres exponentially as \(T \to 0\), which indicates that the system exhibits long-range correlations at low temperatures, and it confirms the presence of a critical point at zero temperature in the 1D Ising model. This behavior of the correlation length is consistent with the absence of a phase transition at finite temperature, as the correlation length remains finite for all \(T > 0\) and only diverges as \(T \to 0\).

The correlation lenght correspond to the distance after which we can consider the spins as independent, since the correlation function decays exponentially with distance \(r\) as \(e^{-r/\xi}\). As the temperature approaches zero, the correlation length diverges, indicating that the spins become correlated over longer distances, which is a signature of critical behavior at zero temperature. This analysis of the two-point correlation function and the correlation length provides insights into the nature of correlations in the 1D Ising model and helps to understand the absence of a phase transition at finite temperature. But since this divergence occur at \(T=0\) then everything is in accord to the fact that there is no phase transition at finite temperature, and the critical point is located at zero temperature.
